\chapter{System as built}
\label{ch:method}

\Cref{ch:methodsmain} gives the command mapping and safety architecture in
summary. This appendix covers what is not repeated there: the package
architecture, the full mode registry, the desk arrangement the laboratory
actually operates, calibration, the operating window, and the verification
audit.

\section{Architecture}

The workspace is six packages with a one-way dependency graph
(\cref{fig:arch}). Direct teleoperation is the baseline condition of every
planned experiment, so if the teleoperation package imported anything from
the autonomy package the comparison would be circular; verified both by
source inspection and by the inverse-kinematics follower still reaching its
running state with the autonomy and perception packages deleted from the
install tree.

\begin{figure}[htbp]
  \centering
  % SYSTEM ARCHITECTURE. Packages, and the one-way dependency arrow that makes
% the experimental comparison non-circular.
\begin{tikzpicture}[x=1mm, y=1mm, font=\scriptsize]

\node[draw, fill=blue!8, minimum width=34mm, minimum height=13mm,
      align=center] (teleop) at (0, 0)
  {\textbf{srl\_teleop}\\ master sensing, IK follower,\\ clutch, e-stop, bridge};
\node[draw, fill=blue!8, minimum width=30mm, minimum height=13mm,
      align=center] (percep) at (46, 0)
  {\textbf{srl\_perception}\\ AprilTag, object pose};
\node[draw, fill=gray!10, minimum width=34mm, minimum height=13mm,
      align=center] (auto) at (23, -22)
  {\textbf{srl\_autonomy}\\ grasps, intent, arbiter};
\node[draw, fill=gray!10, minimum width=34mm, minimum height=13mm,
      align=center] (exp) at (23, -42)
  {\textbf{srl\_experiments}\\ protocols, logging, analysis};
\node[draw, fill=gray!10, minimum width=32mm, minimum height=11mm,
      align=center] (vr) at (-46, -22)
  {\textbf{srl\_vr\_teleop}\\ vendor bridge};
\node[draw, fill=gray!6, minimum width=32mm, minimum height=11mm,
      align=center] (desc) at (-46, 0)
  {\textbf{srl\_description}\\ URDF: arms, mount, wearer};
\node[draw, fill=gray!6, minimum width=32mm, minimum height=11mm,
      align=center] (moveit) at (-46, 16)
  {\textbf{srl\_moveit\_config}\\ TRAC-IK, SRDF, controllers};

\draw[flow] (auto) -- (teleop);
\draw[flow] (auto) -- (percep);
\draw[flow] (exp) -- (auto);
\draw[flow] (exp.west) to[out=180,in=270] (teleop.south west);
\draw[flow] (desc) -- (moveit);
\draw[flow] (moveit.east) to[out=0,in=150] (teleop.north west);

\node[draw, dashed, thick, red!60!black, fill=red!4, text width=62mm,
      align=left] at (0, -62)
  {\textbf{\color{red!60!black}The dependency arrow runs one way, and it is
   load-bearing.} \texttt{srl\_teleop} is the baseline condition of every
   experiment. If it imported anything from \texttt{srl\_autonomy} the
   comparison would be circular. Verified two ways: no such import exists, and
   the IK follower still reaches \emph{Tracking enabled} with the autonomy and
   perception packages physically removed from \texttt{install/}.};
\end{tikzpicture}

  \caption[Package architecture]{Package decomposition and the one-way
  dependency graph.}
  \label{fig:arch}
\end{figure}

\section{Command mapping, in full}

The commanded position is built in spherical form about the master's
shoulder rather than by forward kinematics from seven joint angles, each
component taken from whichever sensor observes it most directly:
\begin{align}
  \phi &= \arcsin\!\left(\hat{a}\cdot\frac{\mathbf{g}}{\lVert\mathbf{g}\rVert}\right)
    &&\text{elevation, from gravity} \label{eq:elev}\\
  \psi &= \theta_1-\theta_1^0 &&\text{azimuth, from the shoulder roll} \label{eq:azim}\\
  r &= \lVert\mathrm{FK}(\theta_1,\theta_2,\theta_3,\theta_4,0,0,0)\rVert &&\text{radius, magnitude only} \label{eq:reach}
\end{align}
with $\hat{a}$ the arm's long axis in the sensor's own frame, captured once
at rest, and $\mathbf{g}$ the measured acceleration. Elevation is
mount-independent and drift-free; radius is dominated by the bend joints and
tolerates a dead roll passed through as zero. The commanded end-effector
position is
\begin{equation}
  \mathbf{p} = \mathbf{p}_{\text{anchor}} + s(\mathbf{d}-\mathbf{d}_{\text{ref}}),
  \label{eq:anchor}
\end{equation}
with $\mathbf{d}$ the master tip in the robot frame and
$\mathbf{d}_{\text{ref}}$ the reference latched at clutch engagement, from
the first \emph{valid} frame after boot at the master's actual pose. An
earlier version initialised the reference to the origin of the master's own
frame, so the first command was a step of the master's entire tip vector
(roughly \SI{0.22}{\metre}) in whatever direction the master happened to be
lying.

\subsection{Collision avoidance in three stages}
\label{sec:clearance}

The wearer is in the collision model: thirteen collision elements were
added, paired with thirty-six exclusions for the mount-adjacent proximal
links (the arm bases are bolted to the harness and permanently adjacent, so
without exclusions every pose is trivially in collision -- everything distal
of the first arm link stays enabled, which is what actually protects the
wearer). Avoidance is graduated: collision-aware inverse kinematics rejects
a configuration; on rejection the redundant degree of freedom is re-seeded
and the solve retried; a geometric clearance floor measured from the
transform tree provides the final margin. The floor is measured against the
\emph{fused} body of \cref{sec:wearer-tracking}, not the semantic collision
description, which excludes forty-four proximal pairs and would report a
pose valid with the upper-arm tube inside the person; and
\texttt{base\_link} sits \SI{0.1610}{\metre} from the torso and is moved by
no joint, so a clearance reading of exactly that value is a measurement of
the mount, not the arm (\cref{ch:workspace}). Driving the commanded pose
directly at the wearer's head and torso stopped the arm at
\SIrange{0.072}{0.079}{\metre} with zero clearance-floor blocks, so the
earlier stages did the work and the floor remained a backstop.

\section{Control modes, in full}

\begin{table}[htbp]
  \centering
  \small
  \caption{Every delivered operating mode. The numbering is 01, 02, 03, 04
  and 06: there is no 05, and the gap is preserved because recorded data
  carries these labels.}
  \label{tab:modes}
  \begin{tabular}{@{}llll@{}}
    \toprule
    Mode & Input & Robot decides & $v_{\max}$ (\si{\radian\per\second}) \\
    \midrule
    01 direct, mannequin & master arm & nothing & 0.60 \\
    02 direct, controllers & tracked controllers & nothing & 0.60 \\
    03 shared autonomy & master, vision & wrist orientation & 0.60 \\
    04 shared autonomy, controllers & controllers, vision & wrist orientation & 0.60 \\
    06 full autonomy & typed or spoken goal & target, approach, grasp, place & 0.15 \\
    \bottomrule
  \end{tabular}
\end{table}

Orientation assist is a stub: it would supply wrist orientation under an
orientation constraint, and the inverse-kinematics plugin in use ignores the
constraint field entirely (270 solver calls, success identical with and
without it). It needs a constraint-aware plugin or explicit yaw sampling.

\subsection{The desk arrangement}
\label{sec:desk}

The laboratory's actual arrangement seats the operator across the room,
facing the wearer, holding the VR controllers as a six-degree-of-freedom
input device and watching the physical robot directly. \textbf{Nobody wears
the headset}; it stands on a shelf as the tracking reference for the
controllers. The alignment between the operator's frame and the robot's is a
measured yaw and can never be a mirror: facing someone and copying their
motion is a reflection (determinant $-1$), which looks correct in position
alone but reverses every commanded orientation. The angle is measured by
having the operator move \SI{0.4}{\metre} toward the robot, with a
perpendicular cross-check; calibration refuses to write a value if the two
motions are not perpendicular and horizontal.

Bringing the WebXR transport up for the first time (2026-08-19) exposed two
defects no unit test would have found: the bridge node had never started at
all, because an attribute named \texttt{clients} shadowed a read-only
property of the base class; and its axis map had been written for the
conventional robotics frame rather than this project's, so a hand moved
\SI{3}{\metre} forward commanded \SI{3}{\metre} to the wearer's right. The
delivered route serves the client page and the \texttt{wss} socket from one
origin so a single certificate exception covers both, needs no \texttt{adb}
or sideloading, and was verified end to end on 2026-08-19.

\section{Calibration}

\begin{figure}[htbp]
  \centering
  % THE CALIBRATION PROCEDURE, IN ORDER.
% The order is not cosmetic: each step consumes the previous step's output,
% and two of the orderings are load-bearing. The notes that say WHY sit
% beside the steps they constrain, brace-linked to them, so a reader can see
% which step each warning is about without a legend.
\begin{tikzpicture}[x=1mm, y=1mm, font=\scriptsize,
  pstep/.style = {proc, text width=44mm, minimum height=11mm},
  why/.style  = {draw, dashed, rounded corners=2pt, red!60!black,
                 fill=red!3, text width=88mm, align=left,
                 font=\scriptsize, inner sep=4pt, anchor=west}]

\node[pstep] (settle) at (0,    0) {\textbf{0. Settle the stream}\\node running, port held,\\IMU past power-on};
\node[pstep] (gyro)   at (0,  -20) {\textbf{1. Gyroscope bias}\\sample \texttt{/master\_arm\_raw\_*},\\never the port};
\node[pstep] (zero)   at (0,  -40) {\textbf{2. Zero reference}\\arm hanging; captures $\hat{\mathbf{a}}$\\and the j1 datum};
\node[pstep] (chan)   at (0,  -60) {\textbf{3. Channel health}\\one verdict per channel over\\\emph{distinct} updates};
\node[pstep] (maxp)   at (0,  -80) {\textbf{4. Range}\\sweep each joint to both\\mechanical stops};
\node[pstep] (btn)    at (0, -100) {\textbf{5. Button map}\\\texttt{check\_buttons} reports,\\never rewrites};
\node[actuator, text width=44mm, minimum height=11mm] (anchor) at (0, -120)
  {\textbf{6. Workspace anchor}\\\texttt{tf2\_echo} at home $\Rightarrow$\\$\mathbf{p}_{\text{HOME}}$ and the anchor quaternion};

\foreach \a/\b in {settle/gyro, gyro/zero, zero/chan, chan/maxp,
                   maxp/btn, btn/anchor}
  {\draw[flow] (\a) -- (\b);}

\node[why] (w0) at (30,  -10)
  {\textbf{Step 0 is not optional.} Re-opening the serial port resets the
   board, and the inertial sensor emits nonsense while it settles. A
   gyroscope bias captured through a script that re-opened the port read
   \SI{4.85}{\degree\per\second} on the left sensor; sampled from the settled
   topic instead, the same sensor read \SI{0.11}{\degree\per\second}. That one
   mistake is the difference between \SI{1492}{\degree\per\minute} and
   \SI{44}{\degree\per\minute} of estimated drift.};
\node[why] (w3) at (30,  -70)
  {\textbf{Step 3 must precede step 4.} A range sweep on an incoherent
   channel measures the fault rather than the joint: the readings span a wide
   interval without forming a trajectory, and range alone calls that healthy.};
\node[why] (w6) at (30, -122)
  {\textbf{Step 6 is derived, not chosen.} The anchor is the end-effector
   pose at the home joint angles, so changing the home pose or the mount
   invalidates it and every clearance figure computed from it.};

\draw[decorate, decoration={brace, amplitude=4pt}, red!60!black]
  (23, -3) -- (23, -17);
\draw[decorate, decoration={brace, amplitude=4pt}, red!60!black]
  (23, -55) -- (23, -85);
\draw[decorate, decoration={brace, amplitude=4pt}, red!60!black]
  (23, -115) -- (23, -126);
\end{tikzpicture}

  \caption[Calibration procedure]{The calibration procedure, in order. A
  calibration routine measures whatever the instrument was actually doing at
  the time, so establishing that the instrument is in a valid state is part
  of the measurement and not a preliminary to it.}
  \label{fig:calibration}
\end{figure}

\section{The operating window}
\label{sec:gui}

Every capability reaches one window: a feature reachable only from a
terminal does not exist for the person running a session with a wearer
standing in the rig. \repo{scripts/verify\_gui\_buttons.py} presses every
control and checks what the press did (257 checks). The first complete run
found four defects: five session controls raised an exception inside a Qt
slot, which the framework swallows, so they did nothing and said nothing;
and the stray-process preflight for the one-stack constraint had never
executed, because its module was never imported.

\begin{figure}[htbp]
  \centering
  \includegraphics[width=0.85\linewidth]{gui/operating_window.png}
  \caption[The operating window]{One command opens it. Left: set-up, run,
  status. Centre: mode registry and per-mode controls. Right: the
  three-dimensional commanded view beside a panel drawn from the arms' own
  reported joint states, with the divergence between them read out beneath
  both, so a mismatch between commanded and reported state is visible
  without reading a log.}
  \label{fig:gui}
\end{figure}

\section{Verification methodology}

The rule -- a surprising failure is evidence about the instrument until the
instrument has been cleared -- and the ten-repeats-over-a-densified-path
practice are stated in \cref{ch:methodsmain}; \cref{ch:instrument} is the
full case-by-case record.
